\chapter{Bring's Surface}
\label{chap:brings-surface}

\section{Bring's surface as a QOU}
\label{sec:brings-qou}

\begin{remark}[Descent to dimension zero]
  \label{rem:descent-to-zero}
  \uses{thm:descartes, def:substrate-parameter}
By the descent theorem (\hyperref[thm:descartes]{Theorem}), it is sufficient
to study dimension-zero manifolds equipped with their bundles $P$ and $E$.
At a point $p\in M_0$, the substrate parameter
$q(p)= 1 + \lvert\langle L_\theta,\, F_A^{+}\rangle\rvert(p)$
(\hyperref[def:substrate-parameter]{Definition}) assigns a specific quantum
constant $q\in[1,\infty)$ determined by the principal bundle $P$ and the
state bundle $E$.  Here $E$ represents the totality of fermionic matter
packed into spacetime at the Big Bang.  The trivial bundle $E=E_0$
corresponds to empty space, in which every point's location is exactly
known; a non-trivial bundle $E$ begins to fill spacetime with fermions,
entering the regime of the quantum Bring surface whose underlying geometry
is non-commutative (\hyperref[sec:quantum-brings]{Section}).
\end{remark}

\begin{example}[Bring's surface unit circle bundle]
  \label{ex:brings-surface}
  \uses{ex:c-quantum-observable-universe}
Let $S$ denote Bring's surface (genus $4$, $\chi = -6$).  The CR geometry
lives on the unit circle bundle $M_S \subset T^*S$: the Liouville
$1$-form restricts to a contact form $\theta$ on $M_S$.  The natural
projection $\pi \colon S \to \mathbb{P}^1(\mathbb{C})$ (degree-$4$
branched covering with $6$ branch points) lifts to a principal
$\mathrm{SU}(2)$-bundle $P \to M_S$, giving an archimedean QOU
$(M_S, \theta, G, P, E)$.
\end{example}

\begin{remark}[Realisations of Bring's surface]
  \label{rem:brings-realisations}
  \uses{def:brings-surface, def:c-quantum-universe, rem:descartes-leptons}
Over an archimedean field $K$ (e.g. $\mathbb{Q}$, $\mathbb{R}$, or
$\mathbb{C}$), Bring's surface is a compact Riemann surface of genus $4$.
The unit circle bundle $M_S$ in $T^*\Sigma_4$ is a CR manifold of
hypersurface type (real dimension $3$), which provides the archimedean
realisation of the quantum observable universe
(\hyperref[def:c-quantum-universe]{Example}).  Physics interpretations --- lepton
generations, chirality, colour --- are further realisations of the
$S_5$-symmetric structure (\hyperref[rem:descartes-leptons]{Remark}).
\end{remark}

\begin{proposition}[Bring's surface is a quantum observable universe]
  \label{prop:brings-is-qou}
  \uses{def:brings-surface, prop:yekutieli-lifting, thm:qu-to-qou}
Bring's surface $\Sigma_4$ (\hyperref[def:brings-surface]{Definition}) gives rise
to a quantum observable universe over $R$ (with $30 \in R^{\times}$):

\begin{enumerate}
  \item \emph{Quantum universe.}  Set $\mathcal{C} =
\end{enumerate}
   \mathrm{Rep}(\mathrm{SU}_{q}(2))$ over $R$,
   $A = \mathcal{O}_q(\mathrm{SU}(2))$ with Frobenius form given by the
   Haar state, and $\omega$ the forgetful functor to $R$-$\mathrm{Mod}$.
   The differential calculus $\Omega^1(A)$ is the Woronowicz
   $3$-dimensional calculus.
\begin{enumerate}
  \item \emph{Observability.}
\end{enumerate}
   $\mathrm{Hom}(\mathbf{1},\, \mathrm{Ad}(A)) = 0$ by the quantum
   Schur lemma: $\mathrm{SU}_q(2)$ has no non-trivial invariants in the
   adjoint representation.
\begin{enumerate}
  \item \emph{Yekutieli lifting.}  The brane tower satisfies the Yekutieli
\end{enumerate}
   lifting condition (\hyperref[prop:yekutieli-lifting]{Proposition}).
\begin{enumerate}
  \item \emph{Euler characteristic.}  By \hyperref[thm:qu-to-qou]{Theorem},
\end{enumerate}
   $\chi(A) = \mathrm{tr}(\mathrm{id}_A) = 2 - 2g = -6$.
\begin{enumerate}
  \item \emph{Geometric descent} (for $n = 1$, $\dim M_S = 3$):
\end{enumerate}
\[
M_S \;\xrightarrow{\;\mathrm{Reeb}\;}\; \Sigma_4
     \;\xrightarrow{\;\pi\;}\; \mathbb{P}(V)
     \;\xrightarrow{\;\mathrm{branch}\;}\; \{6\;\text{points}\}.
\]
   The degree-$4$ covering $\pi$ has $6$ branch points with total
   ramification $R_{\mathrm{tot}} = 14$ (Riemann--Hurwitz), and the
   ramification profile is determined by the $S_5$-symmetry.

---

\emph{Item (1).}
The quantised coordinate ring $A = \mathcal{O}_q(\mathrm{SU}(2))$ is a
Hopf algebra over $R[q, q^{-1}]$ with antipode $S$, counit $\varepsilon$,
and the Haar state $h_q \colon A \to R$ furnishing a non-degenerate
Frobenius form $\lambda = h_q$ (non-degeneracy follows from the
faithfulness of the Haar state for compact quantum groups).  The category
$\mathcal{C} = \mathrm{Rep}(\mathrm{SU}_q(2))$ is $R$-linear, rigid
monoidal (duals given by contragredient representations), and the forgetful
functor $\omega \colon \mathcal{C} \to R\text{-}\mathrm{Mod}$ is exact,
faithful, and monoidal.  The Woronowicz $3D$-calculus
$\Omega^1(A) = \ker(\varepsilon) / \ker(\varepsilon)^2$ is a
$3$-dimensional $A$-bimodule in $\mathcal{C}$, providing the differential
calculus axiom.

\emph{Item (2).}
The adjoint representation $\mathrm{Ad}(A)$ decomposes under
$\mathrm{SU}_q(2)$ into irreducible representations of spin $j \geq 1$
(the spin-$0$ component vanishes because
$\mathrm{Ad} \cong V_1 \oplus V_2 \oplus \cdots$ with no trivial
summand).  Therefore $\mathrm{Hom}(\mathbf{1}, \mathrm{Ad}(A)) = 0$ by
the quantum Schur lemma.

\emph{Item (3).}
By \hyperref[prop:jet-prolongation]{Proposition}, $J^1(A)$ is Frobenius--Hopf, so
the brane tower $\{A_n\}$ is well-defined.  Each transition
$A_n \hookrightarrow A_{n+1}$ gives a central extension of gerbes banded by
the centre $\mathcal{B} = Z(\mathrm{Aut}(A_n))$.  The obstruction class
$\kappa_n \in H^2(\mathcal{S}, \mathcal{B})$ exists by the standard gerbe
classification theorem, verifying the Yekutieli lifting condition
(\hyperref[prop:yekutieli-lifting]{Proposition}).

\emph{Item (4).}
The categorical trace of $\mathrm{id}_A$ in
$\mathcal{C} = \mathrm{Rep}(\mathrm{SU}_q(2))$ is the quantum dimension
$\mathrm{tr}(\mathrm{id}_A) = [N]_q \cdot [N-1]_q / [1]_q$ evaluated at
$N = 2$ (the rank of $\mathrm{SU}(2)$).  At generic $q$, this recovers
the topological Euler characteristic of $\Sigma_4$:
$\chi(A) = 2 - 2g = 2 - 8 = -6$.

\emph{Item (5).}
The Reeb quotient $M_S \to \Sigma_4$ collapses the $\mathrm{U}(1)$-orbits
of the circle bundle (weight-$0$ subalgebra $A_0 = A^{\mathrm{U}(1)}$).
The covering $\pi \colon \Sigma_4 \to \mathbb{P}(V)$ is the map to the
projective line induced by the linear system $|K_{\Sigma_4}|$; it has
degree $4$ and genus $4$ source, so by Riemann--Hurwitz
$R_{\mathrm{tot}} = 2g(\Sigma_4) - 2 - 4(2g(\mathbb{P}^1) - 2) =
6 - 4(-2) = 14$.
The $S_5$-symmetry of $\Sigma_4$ forces the $14$ units of ramification to
be distributed over exactly $6$ branch points (with ramification indices
determined by the cycle types of $S_5$).
\end{proposition}
\begin{proof}
\emph{Item (1).}
The quantised coordinate ring $A = \mathcal{O}_q(\mathrm{SU}(2))$ is a
Hopf algebra over $R[q, q^{-1}]$ with antipode $S$, counit $\varepsilon$,
and the Haar state $h_q \colon A \to R$ furnishing a non-degenerate
Frobenius form $\lambda = h_q$ (non-degeneracy follows from the
faithfulness of the Haar state for compact quantum groups).  The category
$\mathcal{C} = \mathrm{Rep}(\mathrm{SU}_q(2))$ is $R$-linear, rigid
monoidal (duals given by contragredient representations), and the forgetful
functor $\omega \colon \mathcal{C} \to R\text{-}\mathrm{Mod}$ is exact,
faithful, and monoidal.  The Woronowicz $3D$-calculus
$\Omega^1(A) = \ker(\varepsilon) / \ker(\varepsilon)^2$ is a
$3$-dimensional $A$-bimodule in $\mathcal{C}$, providing the differential
calculus axiom.

\emph{Item (2).}
The adjoint representation $\mathrm{Ad}(A)$ decomposes under
$\mathrm{SU}_q(2)$ into irreducible representations of spin $j \geq 1$
(the spin-$0$ component vanishes because
$\mathrm{Ad} \cong V_1 \oplus V_2 \oplus \cdots$ with no trivial
summand).  Therefore $\mathrm{Hom}(\mathbf{1}, \mathrm{Ad}(A)) = 0$ by
the quantum Schur lemma.

\emph{Item (3).}
By \hyperref[prop:jet-prolongation]{Proposition}, $J^1(A)$ is Frobenius--Hopf, so
the brane tower $\{A_n\}$ is well-defined.  Each transition
$A_n \hookrightarrow A_{n+1}$ gives a central extension of gerbes banded by
the centre $\mathcal{B} = Z(\mathrm{Aut}(A_n))$.  The obstruction class
$\kappa_n \in H^2(\mathcal{S}, \mathcal{B})$ exists by the standard gerbe
classification theorem, verifying the Yekutieli lifting condition
(\hyperref[prop:yekutieli-lifting]{Proposition}).

\emph{Item (4).}
The categorical trace of $\mathrm{id}_A$ in
$\mathcal{C} = \mathrm{Rep}(\mathrm{SU}_q(2))$ is the quantum dimension
$\mathrm{tr}(\mathrm{id}_A) = [N]_q \cdot [N-1]_q / [1]_q$ evaluated at
$N = 2$ (the rank of $\mathrm{SU}(2)$).  At generic $q$, this recovers
the topological Euler characteristic of $\Sigma_4$:
$\chi(A) = 2 - 2g = 2 - 8 = -6$.

\emph{Item (5).}
The Reeb quotient $M_S \to \Sigma_4$ collapses the $\mathrm{U}(1)$-orbits
of the circle bundle (weight-$0$ subalgebra $A_0 = A^{\mathrm{U}(1)}$).
The covering $\pi \colon \Sigma_4 \to \mathbb{P}(V)$ is the map to the
projective line induced by the linear system $|K_{\Sigma_4}|$; it has
degree $4$ and genus $4$ source, so by Riemann--Hurwitz
$R_{\mathrm{tot}} = 2g(\Sigma_4) - 2 - 4(2g(\mathbb{P}^1) - 2) =
6 - 4(-2) = 14$.
The $S_5$-symmetry of $\Sigma_4$ forces the $14$ units of ramification to
be distributed over exactly $6$ branch points (with ramification indices
determined by the cycle types of $S_5$).
\end{proof}

\begin{remark}[QOU 4-tuple for Bring's surface]
  \label{rem:brings-qou-4tuple}
  \uses{def:c-quantum-observable-universe, def:brings-surface, def:substrate-parameter, def:su2-generators, def:so3-generators}
We construct the archimedean quantum observable universe
(\hyperref[def:c-quantum-observable-universe]{Example}) explicitly for Bring's
surface.  The five homogeneous coordinates $x_j$, $j\in\{0,\dots,4\}$,
split into \emph{spacetime} and \emph{chiral} degrees of freedom:

\[
\underbrace{x_0,\, x_1}_{\text{spacetime}}
  \qquad
  \underbrace{x_2,\, x_3,\, x_4}_{\text{chiral}}.
\]

The projective ratio $z = x_0/x_1 \in \mathbb{P}^1(\mathbb{C})$ parametrises the
spacetime base.

\textbf{1. CR manifold $M_S$.}
Let $S$ denote Bring's surface (\hyperref[def:brings-surface]{Definition}).
Because $S$ is a Riemann surface (real dimension $2$), it does not itself
carry a contact form.  The CR manifold is the unit circle bundle
$M_S = \{(p,\eta) \in T^*S : \lvert\eta\rvert_{g_S} = 1\}$,
where $g_S$ is the Fubini--Study metric restricted to $S$.  The Liouville
$1$-form $\lambda = \eta\,dz$ on $T^*S$ restricts to a contact form
$\theta = \lambda\big|_{M_S}$, making $M_S$ a $3$-dimensional strictly
pseudoconvex CR manifold.

\textbf{2. Principal bundle $P$.}
The branched covering $\pi\colon S \to \mathbb{P}^1(\mathbb{C})$ (degree $4$, six
branch points) lifts to a principal $\mathrm{SU}(2)$-bundle
$P \to M_S$.  In the spacetime coordinate $z$, the transition functions
of $P$ are the $\mathrm{SU}(2)$ holonomy matrices
$H_1,H_2$ around the branch cuts; in the
chiral coordinates $(x_2,x_3,x_4)$, the $S_3$ permutation of the three
sheets gives the $\mathrm{SO}(3)$ holonomy representation
(\hyperref[def:su2-generators]{Definition}--\hyperref[def:so3-generators]{Definition}).

\textbf{3. State bundle $E$.}
The associated state bundle is $E = P \times_G \mathbb{P}^1(\mathbb{C})$,
whose fibre at each point of $M_S$ is the chiral
$\mathbb{P}^1(\mathbb{C})$ parametrised by the three roots of the elliptic
cubic.

\textbf{4. Connection, Levi form, and substrate parameter.}
The instanton connection $A$ on $P$ is expressed in the resolvent basis
$(w_1,\dots,w_4)$ (\hyperref[rem:lagrange-resolvents]{Remark}) as the flat
$B_3$-monodromy connection whose holonomy around the six branch points is
generated by the trefoil holonomy matrices.  The Levi form on $M_S$ is

\[
L \;=\; d\theta\big|_{H \times H},
  \qquad H = \ker\theta \subset TM_S,
\]

where $\theta = \eta\,dz\big|_{M_S}$ and the horizontal distribution $H$ is
spanned by the lifts of the chiral differentials $dx_2, dx_3, dx_4$ to
$M_S$.  The substrate parameter
(\hyperref[def:substrate-parameter]{Definition}) then reads

\[
q(p) \;=\; 1 + \lvert\langle L_\theta,\, F_A^+ \rangle\rvert(p),
  \qquad p \in M_S,
\]

coupling the chiral curvature to the contact geometry.  At the six branch
points (where chiral sheets merge), $\lVert F_A^+\rVert$ is maximal and
$q \to \infty$ (deep-quantum regime); away from the defects,
$q \to 1$ (classicality).

\textbf{Summary.}
The $4$-tuple $(M_S,\, G,\, P,\, E)$, with $G = \mathrm{SU}(2)$,
constitutes an archimedean quantum observable universe
(\hyperref[ex:brings-surface]{Example}), and every QOU structure---contact
form, Reeb field, Levi form, substrate parameter, $q$-deformed
operators---is determined by the spacetime coordinate $z = x_0/x_1$ and
the chiral coordinates $(x_2, x_3, x_4)$.
\end{remark}

\begin{remark}[QOU axioms instantiated on Bring's surface]
  \label{rem:brings-qou-axiom-map}
  \uses{def:quantum-observable-universe, def:positive-definite-dagger, rem:brings-qou-4tuple, rem:lagrange-resolvents, rem:archimedean-realisations, lem:volume-from-rigidity}
We record how each axiom of a quantum observable universe
(\hyperref[def:quantum-observable-universe]{Definition}) is realised on Bring's
surface.

\textbf{(i) Non-degeneracy.}
The Frobenius form $\lambda = h_q$ is the Haar state on
$A = \mathcal{O}_q(\mathrm{SU}(2))$.  Faithfulness of $h_q$ for compact
quantum groups gives the isomorphism
$A \xrightarrow{\;\sim\;} A^{\vee}$.  The antipode $S$ and $h_q$
together define the dagger $f^{\dagger} = h_q \circ \mu(S(f) \otimes -)$,
which is positive-definite
(\hyperref[def:positive-definite-dagger]{Definition}) because the Haar state
satisfies $h_q(a^* a) > 0$ for $a \neq 0$.

\textbf{(ii) Observability.}
The adjoint representation decomposes as
$\mathrm{Ad}(A) \cong V_1 \oplus V_2 \oplus \cdots$ with no spin-$0$
summand, so $\mathrm{Hom}(\mathbf{1},\, \mathrm{Ad}(A)) = 0$ by the
quantum Schur lemma.  Geometrically, this is the statement that
Bring's surface has no global holomorphic vector field
(it has genus $g = 4 > 1$).

\textbf{(iii) Categorical volume constraint.}
The fundamental object $V$ of
$\mathcal{C} = \mathrm{Rep}(\mathrm{SU}_q(2))$ has categorical
dimension $2$.  In the archimedean realisation, $V$ is the rank-$2$
defining representation carried by the state bundle fibre
$E_p \cong \mathbb{C}^2$ (\hyperref[rem:brings-qou-4tuple]{Remark}, item 3).
The non-degenerate epimorphism $\epsilon \colon \Lambda^2 V \to \mathbf{1}$
is instantiated by the Levi form
$L_\theta = -i\,d\theta|_{H \times H}$
(\hyperref[rem:archimedean-realisations]{Remark}): because $L_\theta$ is
alternating and non-degenerate on the rank-$4$ real distribution
$H = \ker\theta$, it furnishes the symplectic pairing that makes $V$
self-dual and forces $\det(g_V) = 1$, selecting $\mathrm{SU}(2)$ over
$\mathrm{GL}(2)$.  In the spectral coordinates
(\hyperref[rem:lagrange-resolvents]{Remark}), this constraint appears as the
Hermitian quadric $w_1 w_4 + w_2 w_3 = 0$, which is the $\mathrm{SL}(2)$
Plucker relation for $\Lambda^2 V$.  The top-degree extension to the
full chiral fibre is the spectral volume
constraint; see
\hyperref[rem:seiberg-witten]{Remark} for the explicit symplectic
form and its relation to $\epsilon$.
\end{remark}

\section{Lagrange resolvents and braid-group action}
\label{sec:resolvents-braid}

\begin{remark}[Lagrange resolvents and spectral form]
  \label{rem:lagrange-resolvents}
  \uses{def:brings-surface}
Let $\zeta=e^{2\pi i/5}$.  Apply the discrete Fourier transform to the
coordinates $x_j$ of Bring's surface
(\hyperref[def:brings-surface]{Definition}):

\[
w_k \;=\; \frac{1}{\sqrt{5}}\sum_{j=0}^{4}\zeta^{kj}\,x_j,
  \qquad k\in\{0,1,2,3,4\}.
\]

Under this change of basis the three defining power-sum equations become the
following \emph{spectral forms}:

\begin{align}
  w_0 &= 0
    &&\text{(linear projection),}
    \label{eq:resolvent-linear}\\[4pt]
  w_1\,w_4+w_2\,w_3 &= 0
    &&\text{(Hermitian quadric),}
    \label{eq:resolvent-quadric}\\[4pt]
  w_1^2\,w_3+w_4^2\,w_2+w_2^2\,w_1+w_3^2\,w_4 &= 0
    &&\text{(symmetric cubic).}
    \label{eq:resolvent-cubic}
\end{align}

The cyclic permutation $\sigma\colon x_j\mapsto x_{j+1}$ acts by
$\sigma(w_k)=\zeta^{-k}w_k$, so the $C_5$ symmetry is manifest as a phase
rotation $w_k\mapsto\zeta^k\,w_k$.
\end{remark}

\begin{remark}[Braid-group action on the resolvents]
  \label{rem:braid-group-action}
  \uses{rem:brings-curve, rem:lagrange-resolvents, def:braid-group, def:su2-generators, def:so3-generators, def:braid-twist-flip, rem:chirality-trefoil}
The symmetric group $S_5$ acts on Bring's surface
(\hyperref[rem:brings-curve]{Remark}) by permuting
$(x_0,\dots,x_4)$; this is the maximal automorphism group
(order $120$) for genus $4$.  Three distinguished families of $S_3$
subgroups govern the interplay between the base coordinates $x_j$ and the
spectral resolvents $w_k$
(\hyperref[rem:lagrange-resolvents]{Remark}).  We describe each in terms of the
braid group (\hyperref[def:braid-group]{Definition}), the resolvent coordinates,
and the $\mathrm{SU}(2)$/$\mathrm{SO}(3)$ matrices of
\hyperref[def:su2-generators]{Definition}--\hyperref[def:so3-generators]{Definition}.

\textbf{(i) Spacetime action ($\mathbb{P}^1(\mathbb{C})$ base).}
Fix the chiral-like coordinates $x_2,x_3,x_4$ as internal fiber
degrees of freedom.  The constraint
$w_0=0$ forces $x_0+x_1=C$ (constant), so
the spacetime degrees of freedom collapse projectively to the ratio
$z=x_0/x_1\in\mathbb{P}^1(\mathbb{C})$.  The transposition
$x_0\leftrightarrow x_1$ acts as the Mobius involution
$z\mapsto 1/z$, which is a $\pi$-rotation on $\mathbb{P}^1(\mathbb{C})$ and embeds
into $\mathrm{SU}(2)$ as $\beta=\exp(\pi\,e_1)$.
In the resolvent basis this acts as

\begin{equation}\label{eq:spacetime-shift}
  w_k\;\longmapsto\;
  w_k+\tfrac{1}{\sqrt{5}}\,(1-\zeta^k)(x_1-x_0),
  \qquad k=1,\dots,4,
\end{equation}

an affine shift that fixes the chiral modes and translates the
$\mathbb{P}^1(\mathbb{C})$ fibre.

\textbf{(ii) Spectral $S_3$ (trefoil action).}
The frequency-coupled pairs $(w_1,w_4)$ and $(w_2,w_3)$ are permuted by
the twist and flip involutions
(\hyperref[def:braid-twist-flip]{Definition}).
The full
$S_3$ structure emerges on the projective ratios

\[
\eta_1=\frac{w_1}{w_2},\qquad\eta_2=\frac{w_3}{w_4}:
\]

$\sigma_1$ inverts both ratios $(\eta_1\mapsto\eta_1^{-1},\;
\eta_2\mapsto\eta_2^{-1})$, while $\sigma_2$ exchanges them
$(\eta_1\mapsto\eta_2^{-1},\;\eta_2\mapsto\eta_1^{-1})$.  The
$2$-cycle (reflection) swaps the pairs
$(w_1,w_4)\leftrightarrow(w_2,w_3)$; the $3$-cycle (rotation) sends
$(\eta_1,\eta_2)\mapsto(\eta_1,\,\zeta^j\eta_2)$ for $j\in\{1,2,3\}$,
permuting the three complex branches of the cubic
constraint.  The
$\mathrm{SU}(2)$ representation assigns the $3$-cycle to
$\alpha$ and the reflection to
$\beta$; the $\mathrm{SO}(3)$
images act as the holonomy on the spectral
energy distribution, with $R_\alpha$ rotating the time-like resolvent
pair $(\omega_2,\omega_3)$ and $R_\beta$ reflecting the space-like
axis $\omega_1$.

\textbf{(iii) Coordinate $S_3$.}
Fixing any two coordinates (say $x_0,x_1$) as a spacetime base and
permuting the remaining three $(x_2,x_3,x_4)$ gives a stabiliser
$\mathrm{Stab}(\{x_0,x_1\})\cong S_3\subset S_5$; there are
$\binom{5}{2}=10$ such subgroups.  In the resolvent basis a
permutation $\sigma\in S_3$ of $(x_2,x_3,x_4)$ acts as

\begin{equation}\label{eq:coord-s3-shift}
  w_k\;\longmapsto\;
  w_k+\tfrac{1}{\sqrt{5}}\sum_{j=2}^{4}
    \bigl(\zeta^{k\,\sigma(j)}-\zeta^{kj}\bigr)\,x_j,
\end{equation}

so the $w_k$ transform as vector fields over the $(x_0,x_1)$ plane and
the $S_3$ permutes how these fields couple to the base.  The same
$\mathrm{SO}(3)$ representation realises this
action, now playing the role of the holonomy group on the internal
fibre.

\textbf{Geometric $S_3$.}
When Bring's curve is viewed as a branched cover
$\mathcal{S}\to\mathbb{P}^1$, a third family of $S_3$ subgroups arises:
it permutes the three critical values of the covering map.  These three
families --- spectral, coordinate, and geometric --- correspond to the
three copies of $S_3$ inside $S_5$ identified in
\hyperref[rem:chirality-trefoil]{Remark}.
\end{remark}

\section{Color tube coordinates}
\label{sec:color-tube}

% ERROR: block "color-tube-coordinates" not found
\begin{remark}[Aggregate color quantities]
  \label{rem:aggregate-color-quantities}
  \uses{def:color-tube-coordinates, rem:lagrange-resolvents}
Define the following aggregate quantities in color tube coordinates,
with $\zeta = e^{2\pi i/5}$:

\begin{align}
  \widetilde{S} &\;=\; \sigma + r + g + b
    &&\text{(total color deviation),}
    \label{eq:agg-S-tilde}\\[4pt]
  \widetilde{Q} &\;=\; \sigma b + r g
    &&\text{(paired color product),}
    \label{eq:agg-Q-tilde}\\[4pt]
  \widetilde{D} &\;=\; \zeta^{-1}\sigma + \zeta^{-2}r + \zeta^{-3}g + \zeta^{-4}b
    &&\text{(weighted color deviation),}
    \label{eq:agg-D-tilde}\\[4pt]
  \widetilde{C} &\;=\; \sigma^2 g + b^2 r + r^2\sigma + g^2 b
    &&\text{(intrinsic color cubic).}
    \label{eq:agg-C-tilde}
\end{align}

These are related to the unshifted aggregate quantities
$S = \sum_{k=1}^{4}w_k$ and $D = \sum_{k=1}^{4}\zeta^{-k}w_k$ by

\begin{equation}\label{eq:aggregate-shift}
  S \;=\; \widetilde{S} + 4z, \qquad D \;=\; \widetilde{D} - z,
\end{equation}

where the second identity uses $\sum_{k=1}^{4}\zeta^{-k} = -1$.
The inverse Lagrange resolvent transform
(\hyperref[rem:lagrange-resolvents]{Remark}) gives

\begin{equation}\label{eq:z-from-SD}
  z \;=\; \frac{x_0}{x_1} \;=\; \frac{S}{D}.
\end{equation}
\end{remark}

\begin{remark}[Constraint equations in color tube coordinates]
  \label{rem:color-tube-constraints}
  \uses{rem:aggregate-color-quantities, def:color-tube-coordinates}
The Bring surface constraints take the following form in color tube coordinates.

\textbf{Spacetime recovery.}
Substituting into $z = S/D$
(equation [](#eq:z-from-SD)) yields $z(\widetilde{D} - z) = \widetilde{S} + 4z$,
which rearranges to

\begin{equation}\label{eq:spacetime-recovery}
  z^2 - (\widetilde{D} - 4)\,z + \widetilde{S} \;=\; 0.
\end{equation}

Thus the spacetime coordinate $z$ is determined implicitly by the
shifted data $(\widetilde{S}, \widetilde{D})$.

\textbf{Shifted quadric.}
Substituting $w_k = c_k + z$ into the quadric gives

\begin{equation}\label{eq:shifted-quadric}
  \widetilde{Q} + z\,\widetilde{S} + 2z^2 \;=\; 0,
\end{equation}

or equivalently $\sigma b + rg = -z(\sigma + r + g + b) - 2z^2$.

\textbf{Shifted cubic.}
Substituting $w_k = c_k + z$ into the cubic gives

\begin{equation}\label{eq:shifted-cubic}
  \widetilde{C} + z(\widetilde{S}^2 - 2\widetilde{Q})
  + 3z^2\,\widetilde{S} + 4z^3 \;=\; 0.
\end{equation}

Eliminating $\widetilde{Q}$ via the shifted quadric yields the
\emph{reduced cubic}:

\begin{equation}\label{eq:reduced-cubic}
  \widetilde{C} + z\,\widetilde{S}^2 + 5z^2\,\widetilde{S} + 8z^3 \;=\; 0.
\end{equation}

All three equations are polynomial over $\mathbb{Z}$ in $z, r, g, b, \sigma$;
they therefore define affine varieties over $\mathbb{Z}$ and descend
to $\mathbb{F}_\mathfrak{p}$ for every prime $\mathfrak{p} \ne 5$.
\end{remark}

\begin{remark}[Interpretation of the color tube]
  \label{rem:color-tube-interpretation}
  \uses{def:color-tube-coordinates, rem:braid-group-action, rem:brings-qou-4tuple, rem:spectral-dynamics}
The terminology \emph{color tube} reflects two features.

\textbf{Color.}
The coordinates $r, g, b$ label the three chiral degrees of freedom
that correspond, via the $S_3$ subgroups of $S_5$
(\hyperref[rem:braid-group-action]{Remark}(iii)), to the three
colour charges of quantum chromodynamics.  The chiral coordinates
$x_2, x_3, x_4$ of the QOU $4$-tuple
(\hyperref[rem:brings-qou-4tuple]{Remark}) descend to $r = w_2 - z$,
$g = w_3 - z$, $b = w_4 - z$ after shifting by the spacetime ratio $z$.

\textbf{Tube.}
For fixed $z \in \mathbb{P}^1(\mathbb{C})$, the color vector $\vec{c} = (r, g, b)$
is constrained by the shifted quadric and
cubic to lie on a curve inside the
three-dimensional color space $\mathbb{A}^3$.  As $z$ varies over the
spacetime base, this curve sweeps out a tube --- the total space of
the chiral fibration over $\mathbb{P}^1(\mathbb{C})$.

\textbf{Space-like coordinate.}
The shift $\sigma = w_1 - z$ is distinguished: it measures the deviation
of the $C_5$-fundamental resolvent $w_1$ from the spacetime base,
anchoring the $U(1)$ longitudinal backbone
(\hyperref[rem:spectral-dynamics]{Remark}) in the shifted frame.
\end{remark}

\begin{remark}[$C_5$ action in color tube coordinates]
  \label{rem:c5-color-tube}
  \uses{rem:lagrange-resolvents, def:color-tube-coordinates}
The cyclic symmetry $w_k \mapsto \zeta^{-k} w_k$
(\hyperref[rem:lagrange-resolvents]{Remark}) becomes, in color tube coordinates:

\begin{align}
  \sigma &\;\longmapsto\; \zeta^{-1}\sigma + (\zeta^{-1}-1)\,z,
    \label{eq:c5-sigma}\\[4pt]
  r &\;\longmapsto\; \zeta^{-2}r + (\zeta^{-2}-1)\,z,
    \label{eq:c5-r}\\[4pt]
  g &\;\longmapsto\; \zeta^{-3}g + (\zeta^{-3}-1)\,z,
    \label{eq:c5-g}\\[4pt]
  b &\;\longmapsto\; \zeta^{-4}b + (\zeta^{-4}-1)\,z.
    \label{eq:c5-b}
\end{align}

The $C_5$ action is no longer diagonal: each coordinate acquires an
additive correction proportional to $z$, reflecting the cost of
breaking projective homogeneity by subtracting the affine
quantity $z$ from each resolvent.
\end{remark}

\begin{remark}[Intrinsic color cubic: pairing structure]
  \label{rem:intrinsic-color-cubic-pairing}
  \uses{rem:aggregate-color-quantities}
The intrinsic color cubic $\widetilde{C}$ admits a natural decomposition into space-like and chiral sectors:

\begin{equation}\label{eq:color-cubic-decomp}
  \widetilde{C}
  \;=\;
  \underbrace{\sigma(\sigma g + r^2)}_{\text{space--colour coupling}}
  \;+\;
  \underbrace{b(br + g^2)}_{\text{pure colour}}.
\end{equation}

The four terms pair the resolvent indices as $(1,3)$, $(2,1)$, $(4,2)$,
$(3,4)$: each index is squared and contracted with the index two steps
ahead in the cyclic order $1 \to 3 \to 4 \to 2 \to 1$.  This
skip-one pairing respects the conjugation $w_k \leftrightarrow w_{5-k}$
that pairs $(w_1, w_4)$ and $(w_2, w_3)$.
\end{remark}

\section{Chiral fibration and elliptic curves}
\label{sec:chiral-fibration}

\begin{proposition}[Chiral fibre as an elliptic curve]
  \label{prop:chiral-elliptic}
  \uses{rem:braid-group-action, rem:holonomy-trefoil}
Let $z=x_0/x_1\in\mathbb{P}^1(\mathbb{C})$ be the projective spacetime ratio
(\hyperref[rem:braid-group-action]{Remark}(i)).  Substituting the
constraint $x_0+x_1=C$ and eliminating $x_0,x_1$ in favour of $z$,
the quadratic and cubic power-sum
conditions on
the chiral variable $X=(x_2,x_3,x_4)$ reduce at each fixed $z$ to the
cubic

\begin{equation}\label{eq:chiral-elliptic}
  X^3+(z+1)\,X^2+(z^2+z+1)\,X+(z^3+z^2+z+1)=0.
\end{equation}

For generic $z$ this defines an elliptic curve $E_z$ in the chiral
fibre; its $j$-invariant varies as a rational function of $z$.  The
monodromy of $E_z$ around the singular fibres (where the discriminant
of the cubic vanishes) recovers the holonomy
generators $H_1,H_2$ of
\hyperref[rem:holonomy-trefoil]{Remark}.

---

Write $x_0=z\,x_1$ and use $x_0+x_1=C$ to eliminate both spacetime
variables.  The power-sum conditions $\sum x_j^k=0$ for $k=2,3$ then
become polynomial in $z$ and $X=x_2+x_3+x_4$, $X_2=\sum_{i<j}x_ix_j$,
etc.  Reducing by the Newton identities yields the
cubic; its discriminant is a polynomial
in $z$ whose roots give the singular fibres.
\end{proposition}
\begin{proof}
Write $x_0=z\,x_1$ and use $x_0+x_1=C$ to eliminate both spacetime
variables.  The power-sum conditions $\sum x_j^k=0$ for $k=2,3$ then
become polynomial in $z$ and $X=x_2+x_3+x_4$, $X_2=\sum_{i<j}x_ix_j$,
etc.  Reducing by the Newton identities yields the
cubic; its discriminant is a polynomial
in $z$ whose roots give the singular fibres.
\end{proof}

\begin{remark}[Self-duality of the chiral elliptic curve]
  \label{rem:t-duality-elliptic}
  \uses{prop:chiral-elliptic}
The cubic is symmetric under the exchange
$X\leftrightarrow z$: writing

\[
F(X,z)\;=\;X^3+(z+1)\,X^2+(z^2+z+1)\,X+(z^3+z^2+z+1),
\]

one verifies $F(X,z)=F(z,X)$ term by term.  This involution
$\iota\colon(X,z)\mapsto(z,X)$ is a \emph{T-duality}: it identifies
every chiral fibre $E_z$ with its dual spacetime wrapping $E_X$, so
the internal and external degrees of freedom of the Bring surface are
algebraically indistinguishable.

\textbf{Quotient space.}
Passing to the symmetric coordinates $u=X+z$, $v=Xz$, the curve
becomes

\begin{equation}\label{eq:symmetric-quotient}
  v \;=\; \frac{(u+1)(u^2+1)}{2u+1}\,,
\end{equation}

which is rational in $u$.  Hence the quotient
$E_z/\langle\iota\rangle\cong\mathbb{P}^1(\mathbb{C})$: folding the torus along
its $X=z$ diagonal collapses it to a genus-$0$ Riemann sphere.

\textbf{Weierstrass normal form.}
Depressing the cubic via $X=t-\tfrac{z+1}{3}$ eliminates the quadratic
term:

\begin{equation}\label{eq:weierstrass-depressed}
  t^{3}+p(z)\,t+q(z)=0,\qquad
  p=\frac{2z^2+z+2}{3},\quad
  q=\frac{5(4z^3+3z^2+3z+4)}{27}.
\end{equation}

The discriminant $\Delta(z)=-4\,p(z)^3-27\,q(z)^2$ is a polynomial in $z$
whose roots determine the singular fibres of the elliptic
fibration $\mathcal{S}\to\mathbb{P}^1(\mathbb{C})$.  At those roots the
$j$-invariant $j(E_z)=-1728\cdot 64\,p(z)^3/\Delta(z)$ degenerates.
\end{remark}

\begin{remark}[Color tube cubic]
  \label{rem:color-tube-cubic}
  \uses{prop:chiral-elliptic, def:color-tube-coordinates}
At fixed $z \in \mathbb{P}^1(\mathbb{C})$, the three chiral roots of the
elliptic cubic (\hyperref[prop:chiral-elliptic]{Proposition}) each satisfy the cubic in
the color tube coordinate $\mathbf{c} \in \{\mathbf{r}_z, \mathbf{g}_z, \mathbf{b}_z\}$
obtained by substituting $X = \mathbf{c} + z$
into the chiral elliptic equation.  Expanding gives the
\emph{color tube cubic}:

\begin{equation}\label{eq:color-tube-cubic}
  \mathbf{c}^3 + (4z+1)\,\mathbf{c}^2 + (6z^2+3z+1)\,\mathbf{c}
  + (4z^3+3z^2+2z+1) \;=\; 0.
\end{equation}

This polynomial lies in $\mathbb{Z}[z][\mathbf{c}]$ and defines a
variety over $\mathbb{Z}$.  The leading $z$-coefficients $4, 6, 4$
equal $\binom{4}{1}$, $\binom{4}{2}$, $\binom{4}{3}$, and the
constant term satisfies

\[
4z^3 + 3z^2 + 2z + 1 \;=\; \frac{d}{dz}(z^4 + z^3 + z^2 + z).
\]

Vieta's formulas for the three roots give:

\begin{align}
  \mathbf{r}_z + \mathbf{g}_z + \mathbf{b}_z
    &\;=\; -(4z+1),
    \label{eq:vieta-sum}\\[4pt]
  \mathbf{r}_z\,\mathbf{g}_z + \mathbf{r}_z\,\mathbf{b}_z
    + \mathbf{g}_z\,\mathbf{b}_z
    &\;=\; 6z^2+3z+1,
    \label{eq:vieta-pairs}\\[4pt]
  \mathbf{r}_z\,\mathbf{g}_z\,\mathbf{b}_z
    &\;=\; -(4z^3+3z^2+2z+1).
    \label{eq:vieta-product}
\end{align}
\end{remark}

% ERROR: block "color-tube-fibers" not found
% ERROR: block "observed-value" not found
\begin{remark}[T-duality in color tube coordinates]
  \label{rem:t-duality-color-tube}
  \uses{rem:t-duality-elliptic, rem:color-tube-cubic}
The T-duality involution $\iota \colon (X, z) \mapsto (z, X)$
(\hyperref[rem:t-duality-elliptic]{Remark}) acts on color tube
coordinates as follows.  Writing $c = X - z$, the dual point has
base coordinate $z' = X = c + z$ and chiral coordinate $X' = z$,
giving dual color tube coordinate $c' = z - X = -c$.  Hence

\begin{equation}\label{eq:t-duality-color-tube}
  \iota \colon (c,\, z) \;\longmapsto\; (-c,\; c + z).
\end{equation}

The symmetry $F(X,z) = F(z,X)$
(\hyperref[rem:t-duality-elliptic]{Remark}) translates to covariance of
the color tube cubic: the cubic at
parameter $z$ with root $c$ is carried by $\iota$ to the cubic at
parameter $c + z$ with root $-c$.  One verifies this directly from
$F(c+z,\, z) = F(z,\, c+z)$, since $F(z, c+z)$ is the color tube
cubic at parameter $c+z$ evaluated at $-c$.
\end{remark}

\begin{remark}[Sublocus $\widetilde{C}=0$ and colour confinement]
  \label{rem:color-tube-subloci}
  \uses{rem:color-tube-constraints, rem:intrinsic-color-cubic-pairing}
\textbf{Vanishing of the intrinsic color cubic.}
If $\widetilde{C} = 0$, the reduced cubic factors as

\begin{equation}\label{eq:reduced-cubic-factored}
  z\bigl(\widetilde{S}^2 + 5z\,\widetilde{S} + 8z^2\bigr) \;=\; 0,
\end{equation}

so either $z = 0$ (the spacetime origin) or $\widetilde{S}$ satisfies
the quadratic $\widetilde{S}^2 + 5z\,\widetilde{S} + 8z^2 = 0$.
The discriminant of this quadratic (in $\widetilde{S}$) is
$25z^2 - 32z^2 = -7z^2$, which is negative for real $z \ne 0$.
Over $\mathbb{C}$, the two solutions are

\[
\widetilde{S} \;=\; z\,\frac{-5 \pm \sqrt{-7}}{2};
\]

these are non-real for $z \ne 0$, so the locus $\widetilde{C} = 0$
meets the real locus only at $z = 0$.  Over $\mathbb{F}_\mathfrak{p}$ the
existence of a solution depends on whether $-7$ is a quadratic
residue mod $\mathfrak{p}$.

\textbf{Pure colour sublocus.}
The pure colour term $b(br + g^2)$
vanishes on

\[
\{b = 0\} \;\cup\; \{br + g^2 = 0\}.
\]

The condition $b = 0$ means $w_4 = z$: the fourth resolvent
coincides with the spacetime ratio, collapsing one chiral degree of
freedom onto the base.  The condition $br + g^2 = 0$, i.e.
$(w_4 - z)(w_2 - z) + (w_3 - z)^2 = 0$, is a quadratic relation
among the chiral resolvents and $z$.  Either condition can be
interpreted as a \emph{colour-confinement} constraint: one chiral
mode is frozen in register with the spacetime coordinate, and the
remaining two are coupled by a Hermitian-quadric relation in the
shifted frame.
\end{remark}

\section{Bundle structures and smoothness}
\label{sec:bundle-structures}

\begin{remark}[Chiral space as spectral curve and spinor bundle]
  \label{rem:chiral-bundles}
  \uses{prop:chiral-elliptic, rem:lagrange-resolvents, rem:holonomy-trefoil, rem:t-duality-elliptic}
The spacetime--chiral decomposition of Bring's surface maps naturally onto
three classical bundle structures over the $\mathbb{P}^1(\mathbb{C})$ base $z=x_0/x_1$.

\textbf{Cotangent bundle.}
The three-sheeted cubic (\hyperref[prop:chiral-elliptic]{Proposition}) realises the chiral
variables $X=(x_2,x_3,x_4)$ as a \emph{spectral curve}
$\Sigma\subset T^{*}\mathbb{P}^1(\mathbb{C})$: the base coordinate $z$ parametrises
spacetime position, and the three roots $X_1,X_2,X_3$ of
\hyperref[prop:chiral-elliptic]{Proposition} are eigenvalues of a Higgs-like
field $\Phi$ on the base.  The six singular fibres (where the discriminant
$\Delta(z)$ vanishes) are the loci at which eigenvalues coincide and the
differential $dX$ degenerates in the cotangent fibre.

\textbf{Spinor bundle.}
The spectral modes $w_1,\dots,w_4$
(\hyperref[rem:lagrange-resolvents]{Remark}) live in a rank-$2$ complex vector
bundle over $\mathbb{P}^1(\mathbb{C})$.  The chiral doublet
$\Psi_L=(w_1,w_2)^T$ and its conjugate $\Psi_R=(w_4,w_3)^T$ transform
under the $\mathrm{SU}(2)$ holonomy generators
$H_1,H_2$ of
\hyperref[rem:holonomy-trefoil]{Remark}, making $\Psi_L$ a section of the
\emph{spinor bundle} $\mathcal{S}^{+}\to\mathbb{P}^1(\mathbb{C})$.  The $B_3$
monodromy equips this bundle with a flat connection whose holonomy is the
trefoil-knot group.

\textbf{Adjoint bundle.}
The $\mathrm{SO}(3)$ images $R_\alpha,R_\beta$
act on the triple $(\omega_1,\omega_2,\omega_3)$, which therefore spans a
section of the adjoint bundle $\mathrm{ad}\,P\to\mathbb{P}^1(\mathbb{C})$ associated
to the principal $\mathrm{SU}(2)$-bundle $P$.  The T-duality involution
$\iota\colon X\leftrightarrow z$ (\hyperref[rem:t-duality-elliptic]{Remark})
interchanges the fibre and base of the adjoint bundle, exchanging
spectral eigenvalues with spacetime coordinates.
\end{remark}

\begin{remark}[Smoothness of the Bring surface and resolution of spacetime defects]
  \label{rem:brings-smoothness}
  \uses{def:brings-surface, rem:chiral-bundles, def:briggs-klein-chart, rem:lagrange-resolvents}
The Bring surface $\mathcal{S}$ (\hyperref[def:brings-surface]{Definition}) is a
\emph{smooth} algebraic variety: the three power-sum constraints are
transverse in $\mathbb{P}^4$, so their intersection carries no singular
points.  The six topological defects identified in
\hyperref[rem:chiral-bundles]{Remark}---where the discriminant $\Delta(z)$
vanishes and two chiral sheets $X_i,X_j$ collide---are singularities of
the \emph{projection} $\pi\colon\mathcal{S}\to\mathbb{P}^1(\mathbb{C})$, not of the
surface itself.

\textbf{Resolution in the $x_j$ coordinates.}
At a branch point $z=z_c$ where, say, $x_2=x_3$, define the local chiral
coordinate $\delta=x_2-x_3$.  In a neighbourhood of the defect the
spacetime ratio satisfies

\[
z-z_c\;\sim\;\delta^2,
\]

so while the projection map develops a ramification cusp, $\delta$ passes
smoothly through zero, carrying the geometry from one sheet to the other.
The exceptional divisor of the blow-up at each branch point is therefore
already embedded in the chiral coordinates $x_2,x_3,x_4$: no external
resolution is required.

\textbf{Resolution in the $w_k$ coordinates.}
The DFT mixes all five $x_j$ into every resolvent $w_k$
(\hyperref[rem:lagrange-resolvents]{Remark}), so no single $w_k$ vanishes at a
branch point.  The Briggs--Klein chart
(\hyperref[def:briggs-klein-chart]{Definition}) remains a smooth coordinate
system throughout $0<r<\infty$: the radial parameter $r=\lvert w_1\rvert^2$
is a regular function on $\mathcal{S}$, and the $C_5$ orbits
$w_k\mapsto\zeta^k w_k$ foliate the surface without fixed points.  The
infinite torsion and eigenvalue crossings visible in the base
$\mathbb{P}^1(\mathbb{C})$ are artefacts of projecting away the internal degrees of
freedom that the $w_k$ retain.
\end{remark}

% ERROR: block "briggs-klein-chart" not found
\begin{remark}[Radial projection of Bring's surface]
  \label{rem:radial-projection}
  \uses{def:briggs-klein-chart}
Restricting a Briggs--Klein affine chart
(\hyperref[def:briggs-klein-chart]{Definition}) to the interval $r\in(0,1)$ and
projecting Bring's surface onto this radial coordinate yields a fibration
over the base manifold $M$ with trivial bundle $E=E_0$.  At each
$r$, the coordinate $w_1$ is restricted to a circle of radius $\sqrt{r}$ in
the complex plane $\mathbb{C}_{w_1}$.  This circle is the orbit of the $C_5$
symmetry: the rotation $w_1\mapsto\zeta\,w_1$ moves the point along this
circle without changing $r$.
\end{remark}

\begin{remark}[Four-colourability of Bring's surface]
  \label{rem:four-colourability}
  \uses{def:brings-surface}
By the Heawood bound, an orientable surface of genus $g$ has chromatic number
at most $\lfloor(7+\sqrt{1+48g})/2\rfloor$.  For genus $4$ this bound
equals $10$.  Bring's surface, however, admits a proper colouring with only
four colours --- considerably fewer than the maximum permitted by its genus.
This low chromatic number reflects the high degree of symmetry: the $S_5$
automorphism group constrains the combinatorics of any cell decomposition,
forcing a more efficient colouring than a generic genus-$4$ surface would
require.
\end{remark}

\begin{remark}[The four-colour theorem and Bring's surface]
  \label{rem:four-colour-theorem}
  \uses{rem:four-colourability, rem:radial-projection}
The classical four-colour theorem asserts that every planar graph (equivalently,
every graph embeddable in $S^2$) is four-colourable.  For higher-genus
surfaces the chromatic number can be much larger: the Heawood conjecture
(proved by Ringel--Youngs for $g\geq 1$) shows that the maximum chromatic
number grows as $O(\sqrt{g})$.  That Bring's surface attains the planar
bound of four --- despite living at genus $4$ --- is a reflection of the
rigidity imposed by its $S_5$ symmetry.  In the language of the Descartes
universe, the trivial bundle $E=E_0$
(\hyperref[rem:radial-projection]{Remark}) carries no torsion, and the resulting
flatness of the tessellation forces the same chromatic economy as the
sphere.
\end{remark}

\begin{remark}[Cyclotomic reduction and $C_{\mathfrak{p}}$ preservation]
  \label{rem:cyclotomic-reduction}
  \uses{def:brings-surface, rem:lagrange-resolvents}
Let $\mathfrak{p}$ be a prime (not to be confused with a point $p\in M$).
Working over the cyclotomic ring

\begin{equation}\label{eq:cyclotomic-ring}
  \mathbb{Z}_{\zeta_{\mathfrak{p}}}
  \;:=\;\mathbb{Z}[i,\,\zeta_{\mathfrak{p}}],
\end{equation}

where $\zeta_{\mathfrak{p}}=e^{2\pi i/\mathfrak{p}}$ is a primitive
$\mathfrak{p}$-th root of unity, consider the sublocus

\[
S_{\mathfrak{p}}
  \;:=\;
  \Sigma_4\;\cap\;\mathbb{Z}_{\zeta_{\mathfrak{p}}}^{\,5},
\]

that is, the set of points on Bring's surface whose resolvent coordinates
lie in $\mathbb{Z}_{\zeta_{\mathfrak{p}}}$.  The cyclic subgroup
$C_{\mathfrak{p}}\subset S_5$ generated by the $\mathfrak{p}$-cycle preserves
$S_{\mathfrak{p}}$: for any $P\in S_{\mathfrak{p}}$, the image
$\sigma^k(P)$ remains in $\mathbb{Z}_{\zeta_{\mathfrak{p}}}^{\,5}$ because
$\sigma$ acts by multiplication by $\zeta_{\mathfrak{p}}^{-k}$ in the
resolvent basis (\hyperref[rem:lagrange-resolvents]{Remark}), and
$\zeta_{\mathfrak{p}}^{-k}\in\mathbb{Z}_{\zeta_{\mathfrak{p}}}$.  This
arithmetic compatibility between the cyclic symmetry and the cyclotomic ring
underlies the descent of the $C_5$ action to finite characteristic.
\end{remark}

\section{Quaternionic spectral resolvents}
\label{sec:quaternionic-resolvents}

% ERROR: block "quaternionic-spectral-resolvents" not found
\begin{remark}[Algebraic constraints in the $\omega_k$ basis]
  \label{rem:omega-constraints}
  \uses{def:quaternionic-spectral-resolvents, prop:self-dual-curvature}
In the spectral resolvent basis
(\hyperref[def:quaternionic-spectral-resolvents]{Definition}) the
Bring's surface equations take the following form:

\begin{itemize}
  \item \emph{Barycentric:} $\omega_0=0$.
  \item \emph{Quadric ($\mathbb{H}$ condition):}
\end{itemize}
  $\omega_1\bar{\omega}_1+\omega_2\bar{\omega}_3=0$.
  This preserves the horizontal self-duality $F_A|_H={*_H}F_A|_H$
  (\hyperref[prop:self-dual-curvature]{Proposition}) and enforces a
  Minkowski-like balance between the space-like and time-like modes.
\begin{itemize}
  \item \emph{Cubic (phase locking):} a quintic resonance condition that
\end{itemize}
  locks the expansion of $\ln_q$ to $\exp_q$, preserving the $C_5$
  symmetry and the Descartes condition $\chi=-6$.
\end{remark}

\begin{remark}[Spectral dynamics and the $\mathrm{Sp}(1)$ action]
  \label{rem:spectral-dynamics}
  \uses{def:quaternionic-spectral-resolvents, def:quantum-planck-constant, def:q-reeb-flow, rem:q-harmonic-physical}
The growth profiles in
\hyperref[def:quaternionic-spectral-resolvents]{Definition} separate the space-like
and time-like sectors of the Bring's surface:

\begin{itemize}
  \item \emph{$\omega_1$ (space-like):} the factor $1/\ln_q(1+r)$ diverges
\end{itemize}
  as $r\to 0$, so $\omega_1$ achieves infinite density at the base stratum
  $M_0$, characterising the primordial vortex.  This mode anchors the
  $U(1)$ longitudinal backbone of the $\mathrm{SU}(2)$ bundle.
\begin{itemize}
  \item \emph{$\omega_2,\omega_3$ (time-like):} the factor
\end{itemize}
  $\exp_q(r)-1$ vanishes at $r=0$ and grows rapidly for small $r$,
  driven by the $q$-deformed Reeb flow
  (\hyperref[def:q-reeb-flow]{Definition}).  These ladder modes expand at the big
  bang before tapering as the geometric viscosity
  (\hyperref[rem:q-harmonic-physical]{Remark}) takes hold.

To realise the rotation as a function of $r$, embed the radial parameter
into $\mathrm{Sp}(1)$ via the unit quaternion

\begin{equation}\label{eq:sp1-quaternion}
  \mathbf{q}(r)
  \;=\;
  \cos_q\!\theta(r)+\mathbf{u}\,\sin_q\!\theta(r),
  \qquad
  \theta(r)=\tfrac{\pi}{2}\,e^{-r},
\end{equation}

where $\cos_q$ and $\sin_q$ are the $q$-deformed trigonometric functions
(built from $\exp_q$)

\[
\cos_q(x)=\frac{\exp_q(ix)+\exp_q(-ix)}{2},
  \qquad
  \sin_q(x)=\frac{\exp_q(ix)-\exp_q(-ix)}{2i},
\]

and $\mathbf{u}\in\operatorname{Im}\mathbb{H}$ is a fixed unit imaginary
quaternion.  The conjugation
$\boldsymbol{\omega}'=\mathbf{q}(r)\,\boldsymbol{\omega}\,\mathbf{q}(r)^{*}$
gives maximum twist ($\theta=\pi/2$) in the deep-quantum regime $r\to 0$
and classical freeze ($\theta\to 0$) as $r\to\infty$.
\end{remark}

\begin{remark}[Primordial vortex and the deep-quantum limit]
  \label{rem:primordial-vortex}
  \uses{def:quantum-planck-constant, def:q-torsion-leaf, rem:heisenberg-commutator}
As $r\to 0$ the spectral resolvents converge toward the
\emph{primordial vortex}:

\begin{itemize}
  \item The quantum Planck constant
\end{itemize}
  $\hbar_q\to 1$ (\hyperref[def:quantum-planck-constant]{Definition}).
\begin{itemize}
  \item The $q$-torsion leaves
\end{itemize}
  (\hyperref[def:q-torsion-leaf]{Definition}) collapse to the zero-dimensional
  stratum $M_0=\{\pm 1\}$.
\begin{itemize}
  \item The $\omega_0$ coordinate, although identically zero in the
\end{itemize}
  balanced state, formally blows up in the jet prolongation, representing
  the concentration of energy at a singular spatial point.

Space-like energy is trapped because the time-like rotation becomes infinite
in frequency but zero in displacement, effectively turning the
$\mathrm{SU}(2)$ holonomy into a localised topological charge.
\end{remark}

\begin{remark}[Geometric drag and the mass residue]
  \label{rem:geometric-drag}
  \uses{def:substrate-parameter, def:mass-identity}
The Levi form $L$ (\hyperref[def:substrate-parameter]{Definition}) must be
preserved under the $\mathrm{Sp}(1)$ conjugation.  Since

\[
\lVert\boldsymbol{\omega}'\rVert^2
  =\lVert\boldsymbol{\omega}\rVert^2
  =\omega_1\bar{\omega}_1+\omega_2\bar{\omega}_2+\omega_3\bar{\omega}_3,
\]

the quadric constraint remains satisfied and
the horizontal self-duality $F_A|_H={*_H}F_A|_H$ is preserved.  The mass of a stable vortex is
recovered as the residue of the interaction between the growth profiles:

\begin{equation}\label{eq:mass-residue}
  m \;\propto\;
  \oint_{\gamma}
  \lVert\nabla_A q\rVert\cdot
  \lVert\Omega_1\wedge\Omega_{2,3}\rVert,
\end{equation}

where $\Omega_k$ denotes the curvature form restricted to the $\omega_k$
sector.  This heuristic is made rigorous by the nonabelian multiplicative
integral of the mass identity (\hyperref[def:mass-identity]{Definition}).
When $q=1$ (trivial bundle), the $\omega_k$ commute and $m=0$;
when $q\to 0$ (maximal twist), energy is trapped in the $\omega_1$ spatial
filament.
\end{remark}

\section{Geometric resolution and Seiberg--Witten theory}
\label{sec:seiberg-witten}

\begin{remark}[Geometric resolution and the Seiberg--Witten differential]
  \label{rem:seiberg-witten}
  \uses{rem:brings-smoothness, rem:lagrange-resolvents, rem:chiral-bundles, def:quantum-observable-universe, lem:volume-from-rigidity, rem:holonomy-trefoil}
The infinite geometric torsion $T$ and the six topological defects on the
$\mathbb{P}^1(\mathbb{C})$ base $z=x_0/x_1$ are strictly artefacts of
dimensional projection.  The total space, when lifted into the
complete coordinate framework, is a globally smooth complex manifold.

\textbf{1. The blow-up and the exceptional divisor.}
Near a spacetime defect $z_c$ where chiral sheets (e.g. $x_2$ and
$x_3$) merge, the singularity is resolved by blowing up the point into
an exceptional divisor parameterised by the local coordinate
$\delta=x_2-x_3$.  The local geometry satisfies

\[
z-z_c\;\sim\;\delta^2,
\]

proving that the spacetime base "bounces" at the defect: the chiral
coordinate continuously glides through the exceptional divisor and
transitions between sheets without encountering a singularity
(cf. \hyperref[rem:brings-smoothness]{Remark}).

In the spectral domain, the Jacobian determinant of the DFT
(\hyperref[rem:lagrange-resolvents]{Remark}) is a constant
$\lvert\det(J)\rvert=1$, so the global volume form never collapses:

\begin{equation}\label{eq:spectral-volume-form}
  \Omega \;=\; dw_1\wedge dw_2\wedge dw_3\wedge dw_4
  \;=\; \det(J)\;dx_0\wedge\cdots\wedge dx_4\Big|_{w_0=0}.
\end{equation}

This is the \emph{spectral volume constraint}: the top-degree analogue
of the categorical volume constraint
(\hyperref[def:quantum-observable-universe]{Definition}, item (iii)).  The categorical
constraint provides a non-degenerate pairing
$\epsilon \colon \Lambda^2 V \to \mathbf{1}$ on the rank-$2$ fundamental
object; the spectral volume constraint
extends this to the full chiral fibre, ensuring that the $4$-form $\Omega$
is everywhere non-degenerate.  Both enforce $\det = 1$: the categorical
form forces $g_V \in \mathrm{SL}(V)$, while the spectral form forces
$\det(J) \in \mathrm{U}(1)$.

The link is made explicit by the symplectic form.  On the categorical
side, the alternating pairing on $V \cong \mathbb{C}^2$ is
(\hyperref[lem:volume-from-rigidity]{Lemma})

\[
\epsilon(v \wedge w)
  \;=\; v_1 w_2 - v_2 w_1,
  \qquad v, w \in V.
\]

On the Brings surface, the spectral coordinates split into conjugate
pairs $(w_1, w_4)$ and $(w_2, w_3)$ with
$\overline{w_k} = w_{5-k}$.  The quadric
constraint is
$w_1 w_4 + w_2 w_3 = 0$,
which is the Plucker relation $\Lambda^2 V \hookrightarrow
V \otimes V$.  The induced symplectic $2$-form on the chiral fibre is

\begin{equation}\label{eq:chiral-symplectic-form}
  \omega_{\mathrm{ch}}
  \;=\; dw_1 \wedge dw_4 + dw_2 \wedge dw_3,
\end{equation}

and $\Omega = \tfrac{1}{2}\,\omega_{\mathrm{ch}} \wedge
\omega_{\mathrm{ch}}$, so the non-degeneracy of $\Omega$
(i.e. $\det(J) \neq 0$) is equivalent to the non-degeneracy of
$\omega_{\mathrm{ch}}$, which is exactly the fibre-functor image of the
categorical volume form $\epsilon$.

The geometric torsion $T$ is thereby localised not as a point
singularity but as a trapped, finite magnetic flux integrated over
the surface of the exceptional $\mathbb{P}^1(\mathbb{C})$ sphere.

\textbf{2. The Seiberg--Witten differential.}
With the singularities geometrically resolved, the dynamics of the chiral
space are completely parameterised by a meromorphic $1$-form embedded in
the cotangent bundle $T^{\emph{}\mathbb{P}^1(\mathbb{C})$, the }Seiberg--Witten
differential*:

\begin{equation}\label{eq:seiberg-witten-differential}
  \lambda_{\mathrm{SW}} \;=\; X\,dz,
\end{equation}

where $X\in\{x_2,x_3,x_4\}$ is the internal momentum (Higgs-field
eigenvalue) defining the spectral curve
(\hyperref[rem:chiral-bundles]{Remark}), and $z$ is the spacetime position.
The exact physical properties of the topological defects---specifically
the mass/energy spectrum generating the torsion---are rigorously
determined by integrating this differential along the fundamental
non-contractible loops ($A$ and $B$ cycles) of the resolved chiral
torus:

\begin{equation}\label{eq:seiberg-witten-periods}
  a \;=\; \oint_{A}\lambda_{\mathrm{SW}},
  \qquad
  a_D \;=\; \oint_{B}\lambda_{\mathrm{SW}}.
\end{equation}

By lifting the Bring surface into the full spatial and spectral
coordinate spaces, the degenerate branch points resolve completely into
smooth exceptional divisors.  The infinite local torsion is geometrically
captured as trapped flux, revealing that the discrete braid-group
($B_3$) holonomy (\hyperref[rem:holonomy-trefoil]{Remark}) is simply the
continuous Kodaira monodromy of a globally smooth Seiberg--Witten gauge
theory.
\end{remark}

\begin{remark}[Reeb vector field on the exceptional divisor]
  \label{rem:reeb-exceptional-divisor}
  \uses{def:q-reeb-vector-field, def:substrate-parameter, def:q-reeb-flow, rem:brings-smoothness, rem:heisenberg-commutator}
The Reeb vector field $\xi_q$
(\hyperref[def:q-reeb-vector-field]{Definition}) is the dynamical
complement to the static topology of the Bring surface.  While the
surface itself provides the handles (genus $4$), the Reeb vector
provides the "wind" blowing through them, and the quantum torsion $q$
sets the strength of that wind.

\textbf{The classical limit ($q\to 1$): cones and frozen geometry.}
If the substrate parameter $q\to 1$
(\hyperref[def:substrate-parameter]{Definition}), the
torsion coupling vanishes and the phase rotation
$w_k\mapsto\zeta^k\,w_k$ loses its temporal component.  The fibre
over each base point ceases to twist as one moves along a path on
$\mathbb{P}^1(\mathbb{C})$: the chiral tubes are frozen in the configuration
determined by the power-sum equations.  Metrically this is a conical
limit---the curvature concentrates at the blow-up points but no
circulating torsion field smooths it out.  In terms of the Reeb
flow (\hyperref[def:q-reeb-flow]{Definition}), $\xi_q\to\xi_1=\xi$ becomes
the classical Reeb field, and the Heisenberg commutator
$[\hat{X},\hat{P}]_q\to 0$
(\hyperref[rem:heisenberg-commutator]{Remark}): the Bring variety
collapses to a non-rotating algebraic surface.

\textbf{Exceptional divisor as mixing interface.}
Blowing up a spacetime defect $z_c$ where chiral sheets merge
(\hyperref[rem:brings-smoothness]{Remark}) replaces the singular point with
an exceptional divisor $E_c\cong\mathbb{P}^1(\mathbb{C})\cong S^2$.  This
sphere acts as the boundary between the macroscopic coordinates
$(x_0,x_1)$ and the internal chiral modes $(w_1,w_2,w_3)$: the
self-intersections of the chiral fibres are resolved, and the
$C_5$ phase rotation now circulates \emph{around} the exceptional
divisor rather than passing through a singularity.

\textbf{Reeb vector as torsion field.}
On the blown-up surface, the Reeb vector field $\xi_q$ is the
infinitesimal generator of the $U(1)$ action that rotates the chiral
fibre above the exceptional divisor.  Concretely, the
$q$-deformed Reeb flow $\phi_t^q$
(\hyperref[def:q-reeb-flow]{Definition}) restricted to $E_c$ generates
the Hopf fibration $S^3\to S^2\cong E_c$, and the substrate
parameter $q$ controls the pitch of this rotation.  The vector
field $\xi_q$ thereby provides the "time-like" direction within the
chiral space: it is the flow that allows the spectral modes $w_k$ to
execute the spinning phase rotation
$w_k\mapsto\zeta^k\,w_k$.

\textbf{Global picture.}
The Bring surface provides the genus-$4$ topology---four handles
through which the chiral tubes thread.  The Reeb vector field gives
the longitudinal flow along these handles.  The quantum torsion $q$
modulates the pitch of this flow: when $q\to\infty$ (deep-quantum,
Big Bang) the spinning is maximal and the Heisenberg uncertainty
$[\hat{X},\hat{P}]_q\to J$; when $q\to 1$ (heat death) the spinning
ceases, the Reeb field trivialises, and the topology reduces to a
static cone.
\end{remark}

\begin{remark}[Generation-Dependent Seiberg-Witten Differentials]
  \label{rem:generation-dependent-sw}
  \uses{rem:seiberg-witten, rem:color-tube-cubic, rem:brings-smoothness}
The classical Seiberg-Witten integrals over the $A$ and $B$ cycles fail to distinguish lepton generations. In the QOU framework, the vertical degree of the Reeb decomposition lifts the geometry from the base state bundle $E_1$ to the jet prolongation $J^1E_1$. Consequently, the meromorphic 1-form governing the chiral dynamics is strictly generation-dependent.

For the electron (vertical degree 1), the mass is generated by the baseline Seiberg-Witten differential smoothly integrating over the resolved exceptional divisor $E_c$:

\[
a_D^{(e)} = \oint_{E_c} \lambda_{SW}^{(0)} = \oint_{E_c} X dz = \oint_{E_c} (c + z) dz$$

For the muon (vertical degree 2), the gauge field is lifted to the 1-jet bundle, forcing the integration of the momentum gradient. The prolonged differential is $\lambda_{SW}^{(1)} = dX$. Substituting the implicit differentiation of the color tube cubic yields the explicit dual mass integral:

a_D^{(\mu)} = \oint_{E_c} \lambda_{SW}^{(1)} = \oint_{E_c} \frac{4z^3+3z^2+2z+1}{c \left( 3c^2 + 2(4z+1)c + (6z^2+3z+1) \right)} dz$$
\]

At the topological defect $z_c$ where chiral sheets merge, the discriminant vanishes and $\partial F / \partial c \to 0$. This induces a severe pole in the prolonged differential $\lambda_{SW}^{(1)}$. By Cauchy's Residue Theorem, the integral over the exceptional divisor captures the residue of this jet-bundle gradient. This geometrically inflates the trapped quantum dilogarithm volume, producing the mass hierarchy $m_\mu \gg m_e$ purely from the topology of the blowup, without requiring arbitrarily high crossing numbers.
\end{remark}

\begin{proposition}[Abelianisation of the MI integral on exceptional divisors]
  \label{prop:mi-sw-abelianisation}
  \uses{rem:brings-smoothness, def:mass-identity, rem:chiral-bundles, rem:weight-space-decomposition, prop:jet-prolongation, rem:color-tube-cubic}
Let $E_c \cong \mathbb{P}^1(\mathbb{C})$ be the exceptional divisor at a branch
point $z_c$ (\hyperref[rem:brings-smoothness]{Remark}), and let
$(\nabla_A q,\; L \wedge F^+)$ be the connection--curvature pair of the
mass identity (\hyperref[def:mass-identity]{Definition}).  Then for each
vertical degree $n \geq 0$ in the Reeb decomposition, the restriction
of the nonabelian multiplicative integral to $E_c$ abelianises to the
$n$-th prolonged Seiberg--Witten period:

\begin{equation}\label{eq:mi-sw-abelianisation}
  \bigl\lVert\,
  \mathrm{MI}\!\!\int_{E_c}
  \bigl(\nabla_A q,\; L \wedge F^+\bigr)
  \,\bigr\rVert_{J^n}
  \;=\;
  \bigl\lvert a_D^{(n)}\bigr\rvert
  \;=\;
  \biggl\lvert\oint_{E_c}\lambda_{\mathrm{SW}}^{(n)}\biggr\rvert,
\end{equation}

where $\lVert\cdot\rVert_{J^n}$ denotes the Killing-form norm projected
onto the $n$-th jet component, and $\lambda_{\mathrm{SW}}^{(n)}$ is the
$n$-th prolonged Seiberg--Witten differential.

---

The proof proceeds in three steps.

\emph{Step 1: Abelianisation of the connection.}
On the exceptional divisor $E_c \cong \mathbb{P}^1(\mathbb{C})$ the principal
$\mathrm{SU}(2)$-bundle $P$ trivialises (every $\mathrm{SU}(2)$-bundle
over $S^2$ is classified by $\pi_1(\mathrm{SU}(2)) = 0$).  The
connection $\nabla_A$ therefore reduces to a $\mathfrak{su}(2)$-valued
$1$-form $A|_{E_c}$.  By the weight-space decomposition
(\hyperref[rem:weight-space-decomposition]{Remark}), the restriction splits as

\[
A\big|_{E_c}
  \;=\;
  A^0 \otimes H
  \;+\;
  A^+ \otimes E_+
  \;+\;
  A^- \otimes E_-,
\]

where $H, E_\pm$ are the standard generators of $\mathfrak{su}(2)$.
The Cartan component $A^0$ is a $\mathrm{U}(1)$-valued $1$-form on
$E_c$.  Since $E_c \cong \mathbb{P}^1(\mathbb{C})$ the spectral curve
(\hyperref[rem:chiral-bundles]{Remark}) restricts to a single sheet, and the
Higgs-field eigenvalue is the chiral coordinate $X$.  The abelian
component therefore satisfies $A^0 = X\,dz = \lambda_{\mathrm{SW}}$,
recovering the Seiberg--Witten
differential.

\emph{Step 2: Fake-flatness and the $2$-form datum.}
The fake-flatness condition $F_A = B$ identifies
the curvature $2$-form $L \wedge F^+$ on $E_c$ with the
$\mathfrak{su}(2)$-valued $2$-form $B|_{E_c}$.  The simplicity constraint
forces $B|_{E_c} = F_A^+|_{E_c}$, which on the $2$-sphere $E_c$ is
proportional to the volume form $\mathrm{vol}_{E_c}$.  Consequently
the pair $(\nabla_A q,\; L \wedge F^+)$ restricted to $E_c$ is
\emph{abelian}: both components lie in the Cartan subalgebra after
projection onto the weight-$0$ sector, and the nonabelian MI integral
reduces to the ordinary (abelian) multiplicative integral, which equals
the exponential of the classical contour integral.  Taking the norm
yields the absolute value of the period.

\emph{Step 3: Jet prolongation and vertical degree.}
At vertical degree $n = 0$ the Killing-form projection onto $J^0$
recovers the base differential $\lambda_{\mathrm{SW}}^{(0)} = X\,dz$
and the electron period $a_D^{(e)}$.  At vertical degree $n \geq 1$ the
jet prolongation (\hyperref[prop:jet-prolongation]{Proposition}) lifts the
connection to $J^n E$, replacing $X\,dz$ with the $n$-th prolonged
differential $\lambda_{\mathrm{SW}}^{(n)}$: the $1$-jet component is
$dX$ (the implicit derivative of the color tube
cubic), and higher jets follow by induction.
The MI integral at jet level $n$ is the abelian integral of
$\lambda_{\mathrm{SW}}^{(n)}$ over $E_c$ by the same argument as
Steps 1--2 applied to the prolonged connection $A'$ on $J^n E$.
\end{proposition}
\begin{proof}
The proof proceeds in three steps.

\emph{Step 1: Abelianisation of the connection.}
On the exceptional divisor $E_c \cong \mathbb{P}^1(\mathbb{C})$ the principal
$\mathrm{SU}(2)$-bundle $P$ trivialises (every $\mathrm{SU}(2)$-bundle
over $S^2$ is classified by $\pi_1(\mathrm{SU}(2)) = 0$).  The
connection $\nabla_A$ therefore reduces to a $\mathfrak{su}(2)$-valued
$1$-form $A|_{E_c}$.  By the weight-space decomposition
(\hyperref[rem:weight-space-decomposition]{Remark}), the restriction splits as

\[
A\big|_{E_c}
  \;=\;
  A^0 \otimes H
  \;+\;
  A^+ \otimes E_+
  \;+\;
  A^- \otimes E_-,
\]

where $H, E_\pm$ are the standard generators of $\mathfrak{su}(2)$.
The Cartan component $A^0$ is a $\mathrm{U}(1)$-valued $1$-form on
$E_c$.  Since $E_c \cong \mathbb{P}^1(\mathbb{C})$ the spectral curve
(\hyperref[rem:chiral-bundles]{Remark}) restricts to a single sheet, and the
Higgs-field eigenvalue is the chiral coordinate $X$.  The abelian
component therefore satisfies $A^0 = X\,dz = \lambda_{\mathrm{SW}}$,
recovering the Seiberg--Witten
differential.

\emph{Step 2: Fake-flatness and the $2$-form datum.}
The fake-flatness condition $F_A = B$ identifies
the curvature $2$-form $L \wedge F^+$ on $E_c$ with the
$\mathfrak{su}(2)$-valued $2$-form $B|_{E_c}$.  The simplicity constraint
forces $B|_{E_c} = F_A^+|_{E_c}$, which on the $2$-sphere $E_c$ is
proportional to the volume form $\mathrm{vol}_{E_c}$.  Consequently
the pair $(\nabla_A q,\; L \wedge F^+)$ restricted to $E_c$ is
\emph{abelian}: both components lie in the Cartan subalgebra after
projection onto the weight-$0$ sector, and the nonabelian MI integral
reduces to the ordinary (abelian) multiplicative integral, which equals
the exponential of the classical contour integral.  Taking the norm
yields the absolute value of the period.

\emph{Step 3: Jet prolongation and vertical degree.}
At vertical degree $n = 0$ the Killing-form projection onto $J^0$
recovers the base differential $\lambda_{\mathrm{SW}}^{(0)} = X\,dz$
and the electron period $a_D^{(e)}$.  At vertical degree $n \geq 1$ the
jet prolongation (\hyperref[prop:jet-prolongation]{Proposition}) lifts the
connection to $J^n E$, replacing $X\,dz$ with the $n$-th prolonged
differential $\lambda_{\mathrm{SW}}^{(n)}$: the $1$-jet component is
$dX$ (the implicit derivative of the color tube
cubic), and higher jets follow by induction.
The MI integral at jet level $n$ is the abelian integral of
$\lambda_{\mathrm{SW}}^{(n)}$ over $E_c$ by the same argument as
Steps 1--2 applied to the prolonged connection $A'$ on $J^n E$.
\end{proof}

\section{Quantum Brings}
\label{sec:quantum-brings}

This section will develop the quantum version of the commutator relations in the group relations and Bring's formulas. In particular, the classical commutator $[e_a,e_b] = \epsilon_{abc} e_c$ will be replaced with the $q$-deformed commutator from $U_q(\mathfrak{su}(2))$, describing how the Bring surface equations transform under the $q$-deformed action and relating the non-commutative geometry to the quantum substrate parameter $q(p)$.

\section{Other Euler characteristics}
\label{sec:other-euler}

This section will investigate Bring's surface and its analogues for Euler characteristics other than $\chi = -6$, classifying surfaces for $\chi = -2, -4, -8, \ldots$, determining which automorphism groups arise, and connecting to the varying Euler characteristic discussion.
